\subsection{Project Dictionary}

\ChapterContent{
  Here is the Project Dictionary (Definition of Terms) featuring the key technical and operational concepts we have discussed for your capstone.\\

  \textbf{Asset Inventory}. A systematic process of recording and monitoring the physical resources owned by an entity to ensure financial accuracy and operational readiness (Miller \& Zhao, 2023). In this study, it refers to the digital ledger within the system that tracks the lifecycle, location, and status of X University’s equipment.\\

  \textbf{Business Process Reengineering (BPR)}. The fundamental rethinking and radical redesign of business processes to achieve dramatic improvements in critical measures of performance such as cost, quality, and speed (Hammer \& Champy, 2022). In this study, it serves as the development methodology used to overhaul the university’s manual paper-based filing into an automated digital workflow.
  \textbf{Centralized Database}. A data storage architecture where all information is maintained in a single location, accessible by multiple users from different points of entry (Garner \& Smith, 2024). In this study, it refers to the MariaDB repository that consolidates fragmented departmental spreadsheets into one unified source of truth for the university.\\

  \textbf{Functional Suitability}. The degree to which a product or system provides functions that meet stated and implied needs when used under specified conditions (ISO/IEC 25010, 2023). In this study, it is a key evaluation metric used to determine if the software effectively performs all required inventory tasks for the university staff.\\

  \textbf{ISO 25010}. An international standard for systems and software engineering that defines a quality model for evaluating software products across eight primary characteristics (International Organization for Standardization, 2021). In this study, it provides the formal technical framework used to assess the system’s performance, usability, and reliability.\\

  \textbf{Laravel}. An open-source PHP web framework intended for the development of web applications following the model–view–controller (MVC) architectural pattern (Taylor \& Thompson, 2023). In this study, it is the primary backend software tool used to build the system’s logic, security features, and API integrations.\\

  \textbf{MariaDB}. A community-developed, commercially supported fork of the MySQL relational database management system designed for high-performance data processing (Almeida, 2024). In this study, it is the database engine utilized to store and manage the university’s asset records and user credentials.\\

  \textbf{QR Code (Quick Response Code)}. A type of matrix barcode that contains an optically machine-readable optical label attached to an item that records information related to that item (Tanaka \& Lee, 2024). In this study, it refers to the unique, system-generated tags affixed to equipment that, when scanned by an Android device, trigger instant data retrieval.\\

  \textbf{React Native}. An open-source UI software framework used to develop applications for Android, iOS, and Web by enabling developers to use the React framework alongside native platform capabilities (Facebook Open Source, 2022). In this study, it is the tool used to develop the mobile-integrated scanning module for the university’s inventory auditors.\\

  \textbf{Verification Latency}. The measurable delay between the initiation of an audit request and the final confirmation of an asset’s physical presence (Higgins \& O’Neil, 2023). In this study, it refers to the time-gap the researchers aim to reduce by replacing manual serial number typing with instant QR code scanning.\\
}